\section{模型检验：把拟合误差与测厚误差分开}

前章保留硅的完整往返场厚度与碳化硅的两束模型基准，本章分别核对反射率预测表现及厚度随光学条件的变化。

匹配合成数据的厚度平均绝对相对误差为$0.0008\%$（表\ref{tab:04}）；正式谱形分数在问题二为$1.2074$（表\ref{tab:q2-results}）、问题三为$1.0231$（表\ref{tab:08}），均无量纲。前者以合成光谱的已知厚度为参照，后两者以留出反射率为参照，分别评价厚度恢复与曲线预测。
% src:求解/问题1/结果/合成验证.json:主指标.数值
% src:求解/问题2/结果/验证.json:正式主方法分数_连续留段标准化均方根误差
% src:求解/问题3/结果/回炉轮3候选/仲裁闭环/20260910_131501_81623_建模公式/厚度结果.json:分材料验证.硅.逐折汇总.0.平均标准化均方根误差.两束;分材料验证.硅.逐折汇总.0.平均标准化均方根误差.完整往返;分材料验证.硅.逐折汇总.0.平均标准化均方根误差.训练均值;分材料验证.硅.逐折汇总.0.平均标准化均方根误差.二次趋势;分材料验证.硅.逐折汇总.1.平均标准化均方根误差.两束;分材料验证.硅.逐折汇总.1.平均标准化均方根误差.完整往返;分材料验证.硅.逐折汇总.1.平均标准化均方根误差.训练均值;分材料验证.硅.逐折汇总.1.平均标准化均方根误差.二次趋势;分材料验证.碳化硅.逐折汇总.0.平均标准化均方根误差.两束;分材料验证.碳化硅.逐折汇总.0.平均标准化均方根误差.完整往返;分材料验证.碳化硅.逐折汇总.0.平均标准化均方根误差.训练均值;分材料验证.碳化硅.逐折汇总.0.平均标准化均方根误差.二次趋势;分材料验证.碳化硅.逐折汇总.1.平均标准化均方根误差.两束;分材料验证.碳化硅.逐折汇总.1.平均标准化均方根误差.完整往返;分材料验证.碳化硅.逐折汇总.1.平均标准化均方根误差.训练均值;分材料验证.碳化硅.逐折汇总.1.平均标准化均方根误差.二次趋势

曲线预测仍有覆盖缺口：问题二测试覆盖率为$49.06\%$，低于$75\%$校准分位水平；问题三为$73.91\%$，低于$90\%$名义目标。图\ref{fig:q2-coverage}按观测点统计反射率预测区间的实际覆盖，再与各自的校准分位或名义目标比较。
% src:求解/问题2/结果/区间覆盖.json:测试经验覆盖率;校准经验分位点;构造方法
% src:求解/问题3/结果/回炉轮3候选/仲裁闭环/20260910_131501_81623_建模公式/厚度结果.json:预测区间.4.经验覆盖率;预测区间.4.平均区间宽度_比例;预测区间.5.经验覆盖率;预测区间.5.平均区间宽度_比例;预测区间.6.经验覆盖率;预测区间.6.平均区间宽度_比例;预测区间.7.经验覆盖率;预测区间.7.平均区间宽度_比例;预测区间.20.经验覆盖率;预测区间.20.平均区间宽度_比例;预测区间.21.经验覆盖率;预测区间.21.平均区间宽度_比例;预测区间.22.经验覆盖率;预测区间.22.平均区间宽度_比例;预测区间.23.经验覆盖率;预测区间.23.平均区间宽度_比例;预测区间.36.经验覆盖率;预测区间.36.平均区间宽度_比例;预测区间.37.经验覆盖率;预测区间.37.平均区间宽度_比例;预测区间.38.经验覆盖率;预测区间.38.平均区间宽度_比例;预测区间.39.经验覆盖率;预测区间.39.平均区间宽度_比例;预测区间.52.经验覆盖率;预测区间.52.平均区间宽度_比例;预测区间.53.经验覆盖率;预测区间.53.平均区间宽度_比例;预测区间.54.经验覆盖率;预测区间.54.平均区间宽度_比例;预测区间.55.经验覆盖率;预测区间.55.平均区间宽度_比例;预测区间.4.名义覆盖率
% src:求解/问题3/结果/回炉轮3候选/仲裁闭环/20260910_131501_81623_建模公式/厚度结果.json:逐块评分.4.测试点数;逐块评分.5.测试点数;逐块评分.6.测试点数;逐块评分.7.测试点数;逐块评分.20.测试点数;逐块评分.21.测试点数;逐块评分.22.测试点数;逐块评分.23.测试点数;逐块评分.36.测试点数;逐块评分.37.测试点数;逐块评分.38.测试点数;逐块评分.39.测试点数;逐块评分.52.测试点数;逐块评分.53.测试点数;逐块评分.54.测试点数;逐块评分.55.测试点数

光学条件扰动后，问题一固定合成样本的厚度重估为$6.66$--$10.01\,\mu\mathrm{m}$；问题二合并切分、剖面与参数情景，范围为$0.64$--$18.11\,\mu\mathrm{m}$（图\ref{fig:q2-thickness}）。这些跨度记录厚度与光学参数的相互补偿，硅另考察连续区段差异，碳化硅沿用问题二基准。以下两小节先厘清误差参照，再追问厚度为何随重估改变。
% src:求解/问题1/结果/灵敏度.json:条件压力范围_微米
% src:求解/问题2/结果/厚度结果.json:条件区间_微米
% src:求解/问题3/结果/回炉轮3候选/仲裁闭环/20260910_131501_81623_建模公式/厚度结果.json:核心指标.硅折1完整往返条件厚度_um;核心指标.硅折2完整往返条件厚度_um;核心指标.碳化硅正式基准厚度_um
